\chapter{Legionella Testing}

\section*{What Is This Test?}
Legionella testing detects \textit{Legionella pneumophila} and other \textit{Legionella} species, the causative agents of Legionnaires' disease (a severe pneumonia) and Pontiac fever (a self-limited influenza-like illness without pneumonia). \textit{Legionella} are aerobic, fastidious gram-negative rods that thrive in warm water environments (25--45\textdegree{}C) including building water systems, cooling towers, hot tubs, decorative fountains, and hot water tanks. Transmission is via inhalation of aerosolized contaminated water, not person to person. \textit{L. pneumophila} serogroup 1 accounts for approximately 70--90\% of cases \cite{cunha2016legionnaires}. The urinary antigen test detects \textit{L. pneumophila} serogroup 1 lipopolysaccharide in urine with 70--90\% sensitivity and >99\% specificity \cite{shimada2009systematic}. Respiratory culture on buffered charcoal yeast extract (BCYE) agar with antibiotics is the diagnostic gold standard, detecting all species and serogroups, and enabling molecular typing for outbreak investigations. PCR detects \textit{Legionella} DNA in respiratory specimens with high sensitivity. Serology (indirect fluorescent antibody/IFA) requires acute and convalescent sera and is now primarily used for epidemiologic studies.

\section*{Alternative Names}
\begin{itemize}
  \item Legionella Antigen
  \item Legionella Urinary Antigen
  \item Legionella PCR
  \item Legionella Culture
  \item Legionnaires' Disease Testing
  \item Legionella pneumophila Test
  \item Legionella IFA
\end{itemize}
\section*{Principle}
Urinary antigen test: A lateral-flow immunochromatographic assay or enzyme immunoassay (EIA) detects soluble \textit{L. pneumophila} serogroup 1 lipopolysaccharide antigen excreted in urine. Antigen is detectable 1--3 days after symptom onset and persists for days to weeks after treatment. Sensitivity: 70--90\% for serogroup 1; 0\% for non-serogroup 1 species and serogroups. Specificity: >99\%. Culture: Lower respiratory specimens (sputum, BAL) are inoculated onto BCYE agar supplemented with L-cysteine, iron salts, and antibiotics (polymyxin B, cefamandole, vancomycin, anisomycin). Plates are incubated at 35\textdegree{}C in 2.5--5\% CO$_2$ for 3--7 days. Characteristic colonies show ground-glass appearance. Identification by MALDI-TOF, DFA, or latex agglutination. PCR: Real-time PCR targeting \textit{Legionella}-specific 16S rRNA or \textit{mip} (macrophage infectivity potentiator) gene. Sensitivity 33--100\% depending on specimen type, with high specificity. Provides results in 2--4 hours.

\section*{History}
Legionnaires' disease was first recognized in 1976 when an outbreak of severe pneumonia occurred among attendees of the American Legion convention in Philadelphia (182 cases, 29 deaths) \cite{fields2002legionella}. The causative organism, \textit{Legionella pneumophila}, was isolated by Joseph McDade at the CDC in 1977. The urinary antigen test was developed in the 1980s and commercialized in the 1990s. The recognition of hospital-acquired Legionnaires' disease from contaminated water systems in the 1980s led to development of water safety and monitoring guidelines. Outbreaks in subsequent decades (cooling tower outbreaks in Spain 2001, Bronx NY 2015, Flint MI water crisis associated increase 2014--2015) have reinforced the importance of environmental surveillance and rapid diagnostics \cite{phin2014epidemiology}. The IDSA/ATS CAP guidelines recommend Legionella urinary antigen testing and culture for patients with severe CAP, epidemiologic risk factors, or failure to respond to standard CAP therapy \cite{metlay2019diagnosis}.

\section*{Why This Test Is Ordered}
\begin{itemize}
  \item Diagnosis of Legionnaires' disease in patients with severe community-acquired pneumonia (requiring ICU admission) who have one or more of the following: hyponatremia (Na <130 mEq/L), elevated liver transaminases, gastrointestinal symptoms (diarrhea, nausea), relative bradycardia, or failure to respond to beta-lactam antibiotics
  \item Investigation of pneumonia in patients with epidemiologic risk factors: recent travel (hotel, cruise ship), exposure to hot tubs, recent plumbing work, hospitalized >10 days, immunocompromised state (transplant, corticosteroid therapy, TNF-alpha inhibitors, hematologic malignancy)
  \item Testing during outbreak investigations of Legionnaires' disease to identify the environmental source (matching clinical and environmental isolates by molecular typing)
  \item Diagnosis of hospital-acquired pneumonia when \textit{Legionella} is suspected as a nosocomial pathogen, particularly in facilities with known \textit{Legionella} colonization of the water supply
  \item Evaluation of pneumonia not responding to standard CAP therapy (beta-lactam plus macrolide or fluoroquinolone), as \textit{Legionella} requires specific treatment (levofloxacin, azithromycin, or doxycycline)
\end{itemize}

\section*{Primary vs Secondary Test}
Legionella urinary antigen testing is recommended as a primary test in all patients with severe CAP (requiring ICU), per IDSA/ATS 2019 CAP guidelines. It should also be performed in patients with CAP who have epidemiologic risk factors, and in those who fail to respond to standard therapy. Culture is complementary and essential for detecting non-serogroup 1 species and enabling molecular epidemiology. The test is typically ordered in conjunction with blood cultures, sputum culture, and \textit{S. pneumoniae} urinary antigen.

\section*{How This Test Is Performed}
Urinary antigen: A random, clean-catch urine specimen is collected in a sterile cup. Approximately 3--5 mL is required. No preservatives are needed. The specimen may be refrigerated at 2--8\textdegree{}C for up to 24 hours or frozen for longer storage. Concentration of urine by ultrafiltration may improve sensitivity. Respiratory culture: Sputum (expectorated or induced), tracheal aspirate, BAL, or pleural fluid is collected. Organism is fastidious; the laboratory must be notified specifically that \textit{Legionella} culture is requested, as it requires BCYE agar with supplements that are not part of routine sputum culture media. Specimens should be transported to the laboratory within 2 hours. PCR: Lower respiratory specimens (sputum, BAL) collected and transported in the same manner as for culture.

\section*{What Preparation Is Needed}
No special patient preparation is required for urinary antigen testing or respiratory specimen collection. For sputum collection, rinsing the mouth with water before expectoration reduces oral contamination. Antibiotic therapy (especially fluoroquinolones, macrolides) should ideally be withheld until after specimen collection, though the urinary antigen remains positive for days after initiation of effective therapy.

\section*{Contraindications and Precautions}
\begin{itemize}
  \item No contraindications to testing. The major limitation of urinary antigen testing is that it detects only \textit{L. pneumophila} serogroup 1, which accounts for 70--90\% of legionellosis but misses up to 30\% of cases caused by other serogroups and species (\textit{L. micdadei}, \textit{L. bozemanii}, \textit{L. longbeachae}, etc.). Therefore, a negative urinary antigen test does NOT exclude Legionnaires' disease. Culture and/or PCR should be sent on respiratory specimens when clinical suspicion remains despite a negative urinary antigen. The sensitivity is lower in mild disease and may be reduced in urine specimens collected very early (first 24 hours). Concentrated urine may enhance sensitivity.
\end{itemize}
\section*{Risks}
No significant risks from urine collection. Risks related to sputum induction or bronchoscopy for respiratory specimens are rare. The primary clinical risk is failure to diagnose Legionnaires' disease, which carries a mortality of 10--25\% without appropriate treatment, versus 5\% with prompt appropriate antibiotic therapy.

\section*{What Happens After the Test}
Urinary antigen results are available in 15 minutes to 4 hours (depending on assay format). PCR results in 2--4 hours. Culture: colonies typically appear at 3--5 days; negative cultures held for 7--14 days. Positive results should prompt initiation or continuation of therapy with a fluoroquinolone (levofloxacin) or macrolide (azithromycin), which are the recommended first-line agents. Coverage of \textit{Legionella} should also be included in empiric CAP regimens for severe disease regardless of test availability. Public health authorities should be notified of confirmed cases to facilitate outbreak investigation if indicated.

\section*{Understanding Results}

\subsection*{Below Normal}
\begin{itemize}
  \item \textbf{Urinary antigen negative:} Does not exclude \textit{Legionella} infection. Possible explanations: Infection with \textit{L. pneumophila} non-serogroup 1 (~5--20\% of L. pneumophila disease) or non-\textit{pneumophila} species (10--20\% of legionellosis); early specimen collection before antigen reaches detectable levels (less likely beyond 3 days of symptoms); mild disease with low antigen load; or infection in an immunocompromised patient where antigen production or excretion may be altered
  \item \textbf{Culture negative:} May indicate fastidious organism growth failure, prior antibiotic therapy, or inadequate specimen (sputum vs. saliva). Negative culture at 7--14 days, combined with negative PCR, makes Legionnaires' disease unlikely
  \item \textbf{Negative PCR:} Does not exclude legionellosis, particularly if the specimen is from the upper respiratory tract (throat swab sensitivity is <30\%). Lower respiratory specimens are preferred
\end{itemize}

\subsection*{Above Normal}
\begin{itemize}
  \item \textbf{Urinary antigen positive:} Detects \textit{L. pneumophila} serogroup 1 lipopolysaccharide antigen. Diagnostic of \textit{L. pneumophila} serogroup 1 infection (Legionnaires' disease or, rarely, Pontiac fever with proteinuria). Should be correlated with clinical pneumonia (fever, cough, radiographic infiltrate). Positive predictive value >99\% in patients with pneumonia; false positives are extremely rare
  \item \textbf{Culture positive for \textit{Legionella} species:} Definitive diagnosis. Allows species and serogroup identification, which is critical for outbreak source identification (molecular typing of clinical and environmental isolates by pulsed-field gel electrophoresis/PFGE or whole-genome sequencing/WGS). The most common isolates are \textit{L. pneumophila} serogroup 1 (70--90\%), \textit{L. micdadei} (Pittsburgh pneumonia agent), and \textit{L. longbeachae} (associated with compost/potting soil)
  \item \textbf{PCR positive:} Detects \textit{Legionella} DNA. May remain positive after bacterial viability is lost (after antibiotic therapy). Should be correlated with culture and clinical findings
  \item \textbf{Serology (IFA) positive:} Fourfold rise in antibody titer to $\geq$1:128 between acute and convalescent (3--6 weeks) sera is diagnostic. Single titer $\geq$1:256 in a patient with compatible illness is presumptive evidence. Seroconversion may take 3--6 weeks, limiting utility for acute management
\end{itemize}

\section*{Normal Results}
\textbf{Negative} for \textit{Legionella pneumophila} serogroup 1 antigen in urine. \textbf{No Legionella isolated} from culture. \textbf{Legionella DNA not detected} by PCR. Serology: IFA titer <1:64.

\section*{Follow-up Tests}
\begin{itemize}
  \item \textbf{Legionella culture of respiratory specimen:} Essential when urinary antigen is negative but clinical suspicion remains; detects all species and serogroups, enables antimicrobial susceptibility testing (though resistance is rare and standardized methods are limited) and molecular typing
  \item \textbf{Legionella PCR on respiratory specimen (sputum, BAL):} Rapid confirmatory test; may be positive when culture is negative, especially in patients receiving antibiotics
  \item \textbf{Chest radiograph (PA and lateral):} Evaluation for pneumonia; Legionnaires' disease may present with unilateral lobar consolidation, bilateral infiltrates, or pleural effusion; rapid radiographic progression despite antibiotic therapy is a characteristic finding
  \item \textbf{Comprehensive metabolic panel:} Hyponatremia (sodium <130 mEq/L), elevated transaminases, hypophosphatemia, and elevated CPK are characteristic laboratory abnormalities
  \item \textbf{Environmental water testing:} If nosocomial infection is suspected, environmental sampling of the hospital water supply should be performed to identify the source
\end{itemize}

\section*{References}
\bibliographystyle{unsrtnat}
\bibliography{references}

\clearpage
\section*{Regulatory Status and Authorization Date}
This chapter describes a test category, not a specific commercial product. FDA, CDSCO, EU IVDR, and MHRA approval or registration dates therefore cannot be assigned without the assay manufacturer, product name, instrument, and intended use. For laboratory-developed tests, an individual agency approval date may not exist. See the Preface for the interpretation of regulatory status.

\clearpage
